% fig_traj.tex -- uniform-t CFM eval loss trajectory
% Included via % fig_traj.tex -- uniform-t CFM eval loss trajectory
% Included via % fig_traj.tex -- uniform-t CFM eval loss trajectory
% Included via % fig_traj.tex -- uniform-t CFM eval loss trajectory
% Included via \input{fig_traj} inside a figure environment.
\begin{tikzpicture}
\begin{axis}[
  width=\columnwidth,
  height=4.6cm,
  xlabel={Optimiser step},
  ylabel={Eval loss},
  xlabel style={font=\footnotesize},
  ylabel style={font=\footnotesize},
  tick label style={font=\scriptsize},
  xmin=400, xmax=2100,
  ymin=0.15, ymax=0.40,
  xtick={500,1000,1500,2000},
  ytick={0.15, 0.20, 0.25, 0.30, 0.35, 0.40},
  grid=major, grid style={gray!18, dashed},
  legend style={font=\scriptsize, draw=gray!35, fill=white, fill opacity=0.9,
                text opacity=1,
                at={(0.98,0.97)}, anchor=north east, row sep=-2pt},
  legend cell align=left,
  title={\scriptsize MNIST CFM, 2\,k steps, seed 42},
  title style={yshift=-2pt},
]

% apples-to-apples target bands (+0.05 and +0.10 above baseline plateau 0.1734)
\addplot [domain=400:2100, samples=2, dashed, draw=green!55, thick, forget plot]
  {0.2234};
\addplot [domain=400:2100, samples=2, dashed, draw=red!55, thick, forget plot]
  {0.2734};

\node[font=\scriptsize\itshape, anchor=west, text=green!55!black]
  at (axis cs:520, 0.2288) {$+0.05$ target};
\node[font=\scriptsize\itshape, anchor=west, text=red!55!black]
  at (axis cs:520, 0.2790) {$+0.10$ target};

% baseline trajectory
\addplot [color=pfNorm, thick, mark=*, mark size=1.4pt]
  coordinates {(500,0.2000) (1000,0.1817) (1500,0.1757) (2000,0.1734)};
\addlegendentry{Baseline (DiT attention)}

% PhotonFlow A_ref trajectory
\addplot [color=pfPhotonic, thick, mark=square*, mark size=1.3pt]
  coordinates {(500,0.3804) (1000,0.3303) (1500,0.3075) (2000,0.2687)};
\addlegendentry{PhotonFlow ($A_\mathrm{ref}$)}

% final-point gap annotation
\node[font=\scriptsize, anchor=west, text=pfPhotonic!80!black]
  at (axis cs:1520, 0.298) {\textbf{gap $+0.0953$}};
\draw[pharrow, pfPhotonic!80!black]
  (axis cs:1900,0.293) -- (axis cs:1995,0.272);

% "still descending" annotation, placed below legend to avoid overlap
\node[font=\scriptsize\itshape, anchor=east, text=pfPhotonic!80!black]
  at (axis cs:1980, 0.325) {still descending};

\end{axis}
\end{tikzpicture}
 inside a figure environment.
\begin{tikzpicture}
\begin{axis}[
  width=\columnwidth,
  height=4.6cm,
  xlabel={Optimiser step},
  ylabel={Eval loss},
  xlabel style={font=\footnotesize},
  ylabel style={font=\footnotesize},
  tick label style={font=\scriptsize},
  xmin=400, xmax=2100,
  ymin=0.15, ymax=0.40,
  xtick={500,1000,1500,2000},
  ytick={0.15, 0.20, 0.25, 0.30, 0.35, 0.40},
  grid=major, grid style={gray!18, dashed},
  legend style={font=\scriptsize, draw=gray!35, fill=white, fill opacity=0.9,
                text opacity=1,
                at={(0.98,0.97)}, anchor=north east, row sep=-2pt},
  legend cell align=left,
  title={\scriptsize MNIST CFM, 2\,k steps, seed 42},
  title style={yshift=-2pt},
]

% apples-to-apples target bands (+0.05 and +0.10 above baseline plateau 0.1734)
\addplot [domain=400:2100, samples=2, dashed, draw=green!55, thick, forget plot]
  {0.2234};
\addplot [domain=400:2100, samples=2, dashed, draw=red!55, thick, forget plot]
  {0.2734};

\node[font=\scriptsize\itshape, anchor=west, text=green!55!black]
  at (axis cs:520, 0.2288) {$+0.05$ target};
\node[font=\scriptsize\itshape, anchor=west, text=red!55!black]
  at (axis cs:520, 0.2790) {$+0.10$ target};

% baseline trajectory
\addplot [color=pfNorm, thick, mark=*, mark size=1.4pt]
  coordinates {(500,0.2000) (1000,0.1817) (1500,0.1757) (2000,0.1734)};
\addlegendentry{Baseline (DiT attention)}

% PhotonFlow A_ref trajectory
\addplot [color=pfPhotonic, thick, mark=square*, mark size=1.3pt]
  coordinates {(500,0.3804) (1000,0.3303) (1500,0.3075) (2000,0.2687)};
\addlegendentry{PhotonFlow ($A_\mathrm{ref}$)}

% final-point gap annotation
\node[font=\scriptsize, anchor=west, text=pfPhotonic!80!black]
  at (axis cs:1520, 0.298) {\textbf{gap $+0.0953$}};
\draw[pharrow, pfPhotonic!80!black]
  (axis cs:1900,0.293) -- (axis cs:1995,0.272);

% "still descending" annotation, placed below legend to avoid overlap
\node[font=\scriptsize\itshape, anchor=east, text=pfPhotonic!80!black]
  at (axis cs:1980, 0.325) {still descending};

\end{axis}
\end{tikzpicture}
 inside a figure environment.
\begin{tikzpicture}
\begin{axis}[
  width=\columnwidth,
  height=4.6cm,
  xlabel={Optimiser step},
  ylabel={Eval loss},
  xlabel style={font=\footnotesize},
  ylabel style={font=\footnotesize},
  tick label style={font=\scriptsize},
  xmin=400, xmax=2100,
  ymin=0.15, ymax=0.40,
  xtick={500,1000,1500,2000},
  ytick={0.15, 0.20, 0.25, 0.30, 0.35, 0.40},
  grid=major, grid style={gray!18, dashed},
  legend style={font=\scriptsize, draw=gray!35, fill=white, fill opacity=0.9,
                text opacity=1,
                at={(0.98,0.97)}, anchor=north east, row sep=-2pt},
  legend cell align=left,
  title={\scriptsize MNIST CFM, 2\,k steps, seed 42},
  title style={yshift=-2pt},
]

% apples-to-apples target bands (+0.05 and +0.10 above baseline plateau 0.1734)
\addplot [domain=400:2100, samples=2, dashed, draw=green!55, thick, forget plot]
  {0.2234};
\addplot [domain=400:2100, samples=2, dashed, draw=red!55, thick, forget plot]
  {0.2734};

\node[font=\scriptsize\itshape, anchor=west, text=green!55!black]
  at (axis cs:520, 0.2288) {$+0.05$ target};
\node[font=\scriptsize\itshape, anchor=west, text=red!55!black]
  at (axis cs:520, 0.2790) {$+0.10$ target};

% baseline trajectory
\addplot [color=pfNorm, thick, mark=*, mark size=1.4pt]
  coordinates {(500,0.2000) (1000,0.1817) (1500,0.1757) (2000,0.1734)};
\addlegendentry{Baseline (DiT attention)}

% PhotonFlow A_ref trajectory
\addplot [color=pfPhotonic, thick, mark=square*, mark size=1.3pt]
  coordinates {(500,0.3804) (1000,0.3303) (1500,0.3075) (2000,0.2687)};
\addlegendentry{PhotonFlow ($A_\mathrm{ref}$)}

% final-point gap annotation
\node[font=\scriptsize, anchor=west, text=pfPhotonic!80!black]
  at (axis cs:1520, 0.298) {\textbf{gap $+0.0953$}};
\draw[pharrow, pfPhotonic!80!black]
  (axis cs:1900,0.293) -- (axis cs:1995,0.272);

% "still descending" annotation, placed below legend to avoid overlap
\node[font=\scriptsize\itshape, anchor=east, text=pfPhotonic!80!black]
  at (axis cs:1980, 0.325) {still descending};

\end{axis}
\end{tikzpicture}
 inside a figure environment.
\begin{tikzpicture}
\begin{axis}[
  width=\columnwidth,
  height=4.6cm,
  xlabel={Optimiser step},
  ylabel={Eval loss},
  xlabel style={font=\footnotesize},
  ylabel style={font=\footnotesize},
  tick label style={font=\scriptsize},
  xmin=400, xmax=2100,
  ymin=0.15, ymax=0.40,
  xtick={500,1000,1500,2000},
  ytick={0.15, 0.20, 0.25, 0.30, 0.35, 0.40},
  grid=major, grid style={gray!18, dashed},
  legend style={font=\scriptsize, draw=gray!35, fill=white, fill opacity=0.9,
                text opacity=1,
                at={(0.98,0.97)}, anchor=north east, row sep=-2pt},
  legend cell align=left,
  title={\scriptsize MNIST CFM, 2\,k steps, seed 42},
  title style={yshift=-2pt},
]

% apples-to-apples target bands (+0.05 and +0.10 above baseline plateau 0.1734)
\addplot [domain=400:2100, samples=2, dashed, draw=green!55, thick, forget plot]
  {0.2234};
\addplot [domain=400:2100, samples=2, dashed, draw=red!55, thick, forget plot]
  {0.2734};

\node[font=\scriptsize\itshape, anchor=west, text=green!55!black]
  at (axis cs:520, 0.2288) {$+0.05$ target};
\node[font=\scriptsize\itshape, anchor=west, text=red!55!black]
  at (axis cs:520, 0.2790) {$+0.10$ target};

% baseline trajectory
\addplot [color=pfNorm, thick, mark=*, mark size=1.4pt]
  coordinates {(500,0.2000) (1000,0.1817) (1500,0.1757) (2000,0.1734)};
\addlegendentry{Baseline (DiT attention)}

% PhotonFlow A_ref trajectory
\addplot [color=pfPhotonic, thick, mark=square*, mark size=1.3pt]
  coordinates {(500,0.3804) (1000,0.3303) (1500,0.3075) (2000,0.2687)};
\addlegendentry{PhotonFlow ($A_\mathrm{ref}$)}

% final-point gap annotation
\node[font=\scriptsize, anchor=west, text=pfPhotonic!80!black]
  at (axis cs:1520, 0.298) {\textbf{gap $+0.0953$}};
\draw[pharrow, pfPhotonic!80!black]
  (axis cs:1900,0.293) -- (axis cs:1995,0.272);

% "still descending" annotation, placed below legend to avoid overlap
\node[font=\scriptsize\itshape, anchor=east, text=pfPhotonic!80!black]
  at (axis cs:1980, 0.325) {still descending};

\end{axis}
\end{tikzpicture}
